\section{Related Work and Problem Formulation}
\label{sec:sota}

\subsection{Related Work and State-of-the-Art Taxonomy}
\label{sec:related_work}

Gradient-boosted decision trees (GBDTs), including XGBoost \cite{chen2016xgboost}, LightGBM \cite{ke2017lightgbm} and CatBoost \cite{prokhorenkova2018catboost}, are the standard choice for flow-based network intrusion detection on benchmarks such as CSE-CIC-IDS2018 \cite{sharafaldin2018toward}, UNSW-NB15 \cite{moustafa2015unsw} and CIDDS-001 \cite{ring2017flow}. Comprehensive empirical studies support this choice \cite{dang2022using, dang2023multilayer, apruzzese2022modeling}.

Early adversarial attacks on NIDS were transplanted from image classification. Szegedy et al.~\cite{szegedy2014intriguing}, Goodfellow et al.~\cite{goodfellow2015explaining} and Papernot et al.~\cite{papernot2016limitations} showed that gradient-based perturbations fool deep networks, and optimizers such as the Carlini-Wagner $L_2$ attack \cite{carlini2017towards} were then applied to tabular flow features. These studies reported high success rates, but the perturbed vectors were not network flows: they assign fractional values to packet counts, modify immutable fields such as destination port, and break the conservation between byte totals and packet counts. The problem-space view corrects this. Pierazzi et al.~\cite{pierazzi2020intrusion} showed that an adversarial perturbation is meaningful only when it is realized as a transformation of the physical object, such as a binary or a packet sequence, and that the transformation must preserve functionality and survive the feature extractor. Verkerken et al.~\cite{verkerken2026problem} apply this view to defining host-space adversarial perturbations against NIDS. Surveys and audits of this literature~\cite{apruzzese2022modeling, apruzzese2023real, he2023adversarial, vitorino2023sok} argue that many feature-space attacks assume threat models that cannot be realized with valid packets under RFC protocol specifications, so their success rates overstate the risk once evaluated against true packet parsers and middleboxes. Constrained tabular attacks address feasibility only in part. CAA \cite{simonetto2024constrained} combines a gradient-based attack (CAPGD) with MOEVA \cite{ghamizi2020moeva}, a multi-objective genetic search whose constraint penalties are heuristic and whose population search receives no directional guidance. TabularBench \cite{simonetto2024tabularbench} standardizes such evaluations across deep and tree models. Packet-level tools avoid these problems by operating on real traffic. Traffic Manipulator \cite{han2021traffic} uses particle swarm optimization to mutate PCAP files by perturbing inter-packet intervals, adding payload bytes and injecting chaff packets. IDSGAN \cite{lin2022idsgan} generates adversarial flow records with generative adversarial networks, but its feasibility is only approximated in feature space.

Recent work makes the realism requirement concrete for NIDS. Merzouk et al.~\cite{merzouk2022practicality} show that many published evasion attacks cannot be realized on live traffic, Sheatsley et al.~\cite{sheatsley2022adversarial} and Debicha et al.~\cite{debicha2023advbot, debicha2026targeted} craft domain-constrained or botnet flows that remain valid, Zhang et al.~\cite{zhang2025explainable} build explainable and transferable attacks, and El Shehaby et al.~\cite{elshehaby2026perturbability} score which flow features an attacker can perturb, the mutable/immutable split that FFE builds into its schema. Learned generators~\cite{lin2022idsgan, zolbayar2022nidsgan} and grey-box frameworks~\cite{he2022liuer} trade query cost for a training phase. For tabular data in general, CAPGD and CAA~\cite{simonetto2024constrained}, cost-aware attacks~\cite{kireev2023cost} and recent benchmarks and surveys~\cite{simonetto2024tabularbench, he2026tabattackbench, dyrmishi2025insights, vitorino2024benchmark} show that constraint handling dominates reported success rates. The tree-ensemble literature moved toward faster and deployable attacks~\cite{cascioli2024faster, pashamokhtari2022adiotack}, evasion detection~\cite{devos2023detecting}, and robust training~\cite{vos2021efficient, vos2022robust, calzavara2021feature, calzavara2023verifiable, calzavara2025verifiableboosted, lindqvist2025adversarial}. None of these exploits split geometry while keeping coupled flow invariants exact.

Evaluation practice has also been questioned: CIC-IDS2017 and CSE-CIC-IDS2018 contain labeling errors and duplicated records~\cite{engelen2021troubleshooting, lanvin2023errors, liu2022error}, datasets differ widely in difficulty~\cite{flood2024spectral, layeghy2024benchmarking}, and detectors generalize poorly across datasets~\cite{cantone2024crossdataset, guo2026knowml}. We therefore evaluate on complete, deduplicated captures (Section~\ref{sec:setup}).

Exact methods exploit the piecewise-constant structure of trees. Kantchelian et al.~\cite{kantchelian2016evasion} cast minimal evasion against tree ensembles as a mixed-integer linear program. The formulation is exact but NP-hard in general and accepts only linear constraints, so it cannot express the nonlinear rate relations between bytes and duration that flow records obey. SAT and MaxSAT encodings of tree ensembles \cite{ignatiev2022formal, izza2022tackling} provide formal explanations and robustness checks, and LT-Attack \cite{zhang2020efficient} attacks tree ensembles through continuous leaf approximations. The MILP and LT-Attack operate in unconstrained feature space and do not enforce coupled network invariants. Robust training of tree ensembles \cite{chen2019robust, andriushchenko2019provably, calzavara2020treant} and adversarial retraining of neural NIDS \cite{hashemi2020enhancing} improve robustness within the same feature-space threat model. Black-box query methods such as ZOO \cite{chen2017zoo} and HopSkipJump \cite{chen2020hopskipjump} estimate the decision boundary by sampling the classifier's output, and population or swarm searches mutate parameters without knowing which splits are active.

Table~\ref{tab:sota_taxonomy_expanded} summarizes these paradigms and the gap they leave. Methods that exploit the structure of tree ensembles or work in unconstrained feature space do not produce flows that survive packet parsing, while methods that preserve physical feasibility either treat the detector as a black box and need hundreds of queries per flow, or train a surrogate and transfer, which is cheap per flow but depends on how well the surrogate mimics the target. Flow-Field Evasion (FFE) addresses this gap by steering flows along infinitesimal packet-level generators that keep every step on the flow manifold, and by using the split geometry of the trees to choose the direction.

\begin{table}[!htbp]
\caption{Representative evasion, verification and defense paradigms for tabular and network telemetry. Physical feasibility reports the empirical PFR of this benchmark where it was evaluated (Section~\ref{sec:results}).}\label{tab:sota_taxonomy_expanded}
\centering
\linespread{1}\selectfont
\footnotesize
\setlength{\tabcolsep}{4pt}
\begin{tabular}{@{}p{2.0cm}p{1.6cm}p{2.4cm}p{1.9cm}p{2.1cm}p{2.1cm}@{}}
\toprule
\textbf{Method} & \textbf{Target model} & \textbf{Engine} & \textbf{Queries} & \textbf{Coupled invariants} & \textbf{Physical feasibility} \\
\midrule
Kantchelian et al.~\cite{kantchelian2016evasion} & GBDT, random forest & MILP (worst case $2^{\mathcal{O}(DT)}$) & 0 (white-box) & Unsupported (linear only) & Unconstrained \\
Zhang et al.~\cite{zhang2020efficient} (LT-Attack) & Tree ensembles & Continuous leaf approximation & $O(N_{\text{iter}})$ & Not modeled & Unconstrained \\
Chen et al. (ZOO)~\cite{chen2017zoo} & Black-box classifiers & Zeroth-order coordinate descent & 76--137 (mean) & Not modeled & 14.9--35.6\% (empirical) \\
Chen et al. (HSJ)~\cite{chen2020hopskipjump} & Black-box decision & Boundary binary walk & 47 (mean) & Not modeled & 0.0\% (empirical) \\
Lin et al. (IDSGAN)~\cite{lin2022idsgan}; Zolbayar et al.~\cite{zolbayar2022nidsgan} & Black-box NIDS & Learned generator on a surrogate & 2--5 after training & Repaired after generation & 100.0\% (after repair) \\
Ghamizi et al. (MOEVA)~\cite{ghamizi2020moeva} & Constrained tabular & Multi-objective genetic search & 250--434 (mean) & Heuristic penalty & 48.5--67.5\% (empirical) \\
Simonetto et al. (CAA/CAPGD)~\cite{simonetto2024constrained} & Tabular models & Constrained APGD (surrogate) with MOEVA & 1--15 (CAPGD, mean) & Repair operator & 100.0\% (after repair) \\
Han et al. (Traffic Manipulator)~\cite{han2021traffic} & Black-box NIDS & PSO on packet knobs & 161--262 (mean) & Invariant projection & 100.0\% (projected) \\
\midrule
\textbf{Flow-Field Evasion (ours)} & \textbf{XGBoost, LightGBM} & \textbf{Lie vector field steering} & \textbf{$\le 4M_{\max}+1$ probes + TSBS (mean 13.0--91.4)} & \textbf{Exact closed-form} & \textbf{100.0\% (schema)} \\
\bottomrule
\end{tabular}

\smallskip
\footnotesize{Note: $T$ trees, $D$ depth, $M_{\max}=20$, $K=4$ generators. Empirical query counts and PFR values are ranges over the three full-scale benchmarks against XGBoost (Table~\ref{tab:master_evasion_benchmark}). LT-Attack is cited as white-box context and was not run; CAPGD and the NIDSGAN-style generator are our re-implementations (Section~\ref{sec:setup}).}
\end{table}
